\documentclass[11pt,a4paper]{report}

% ===================== PACKAGES =====================
\usepackage[utf8]{inputenc}
\usepackage[T1]{fontenc}
\usepackage[french]{babel}
\usepackage{lmodern}
\usepackage[margin=2.5cm]{geometry}
\usepackage{xcolor}
\usepackage{tikz}
\usetikzlibrary{arrows.meta,positioning,shapes.multipart,fit,backgrounds,calc,shadows}
\usepackage{booktabs}
\usepackage{longtable}
\usepackage{tabularx}
\usepackage{array}
\usepackage{multirow}
\usepackage{colortbl}
\usepackage{enumitem}
\usepackage{titlesec}
\usepackage{fancyhdr}
\usepackage{rotating}
\usepackage{graphicx}
\usepackage{adjustbox}
\usepackage[skins,breakable]{tcolorbox}
\usepackage{hyperref}
\usepackage{parskip}

% ===================== CHAMPS À COMPLÉTER =====================
% Renseignez ces champs avant compilation (informations académiques).
\newcommand{\ecoleNom}{École Marocaine des Sciences de l'Ingénieur}
\newcommand{\ecoleSigle}{EMSI}
\newcommand{\villeEcole}{Casablanca --- Campus Oranger}
\newcommand{\filiereEtud}{Développement des Systèmes d'Information (DSI) --- 4\textsuperscript{e} année}
\newcommand{\anneeUniv}{2025 -- 2026}
\newcommand{\nomEtudiant}{Marouane Essaadouni}
\newcommand{\nomEncadrant}{Asmaa Roudane}
\newcommand{\typeProjet}{Stage de fin d'année}

% ===================== COULEURS =====================
\definecolor{km0blue}{RGB}{22,52,89}
\definecolor{km0orange}{RGB}{224,122,29}
\definecolor{km0green}{RGB}{45,134,89}
\definecolor{km0red}{RGB}{178,34,34}
\definecolor{km0gray}{RGB}{240,242,245}
\definecolor{km0graytext}{RGB}{90,98,110}

% ===================== TITRES =====================
\titleformat{\chapter}[hang]{\normalfont\Huge\bfseries\color{km0blue}}{\thechapter.}{20pt}{\Huge}
\titleformat{\section}{\normalfont\Large\bfseries\color{km0blue}}{\thesection}{10pt}{}
\titleformat{\subsection}{\normalfont\large\bfseries\color{km0blue}}{\thesubsection}{8pt}{}

% ===================== EN-TETES =====================
\pagestyle{fancy}
\fancyhf{}
\renewcommand{\headrulewidth}{0.4pt}
\fancyhead[L]{\small\textcolor{km0graytext}{KM0 --- Cahier des charges}}
\fancyhead[R]{\small\textcolor{km0graytext}{\leftmark}}
\fancyfoot[C]{\small\thepage}

% ===================== HYPERREF =====================
\hypersetup{
    colorlinks=true,
    linkcolor=km0blue,
    citecolor=km0blue,
    urlcolor=km0blue,
    pdftitle={KM0 - Cahier des charges},
    pdfauthor={Garage KM0}
}

% ===================== ENVIRONNEMENTS UTILES =====================
\newtcolorbox{regle}[2]{
  enhanced, breakable,
  colback=km0blue!4, colframe=km0blue,
  colbacktitle=km0blue, coltitle=white,
  fonttitle=\bfseries,
  title={Règle #1 --- #2},
  boxrule=0.6pt, arc=2pt, top=5pt, bottom=5pt, left=7pt, right=7pt,
  before skip=8pt, after skip=8pt
}

\newcommand{\etat}[3]{%
  % #1 = texte, #2 = couleur, #3 = position tikz optionnelle
  \tikz[baseline=-0.6ex]\node[rectangle,rounded corners=2pt,fill=#2!15,draw=#2,thick,inner sep=3pt]{\textcolor{#2!70!black}{\textbf{#1}}};%
}

% ===================== DOCUMENT =====================
\begin{document}

% --------------------------------------------------------------
% PAGE DE GARDE
% --------------------------------------------------------------
\begin{titlepage}
\centering

% ---- Bandeau supérieur : logo ----
% Placez le fichier logo.png à la racine du projet Overleaf (même dossier que ce .tex).
\includegraphics[width=4.6cm]{logo.png}

\vspace{0.6cm}
{\large\bfseries\color{km0blue} \ecoleNom}\\[0.1cm]
{\normalsize\color{km0graytext} \ecoleSigle\ --- \villeEcole}\\[0.15cm]
{\small\color{km0graytext} Filière~: \filiereEtud}\\[0.1cm]
{\small\color{km0graytext}\bfseries \typeProjet}

\vspace{0.4cm}
\begin{tikzpicture}
\draw[km0blue, line width=1pt] (-7,0) -- (7,0);
\end{tikzpicture}

\vspace{2cm}
{\Huge\bfseries\color{km0blue} CAHIER DES CHARGES}\\[1.2cm]

{\huge\bfseries\color{km0blue} KM0}\\[0.25cm]
{\large\color{km0graytext} Système de gestion de garage automobile}

\vspace{0.3cm}
\begin{tikzpicture}
\draw[km0orange, line width=1.4pt] (-1.8,0) -- (1.8,0);
\end{tikzpicture}

\vspace{0.4cm}
{\color{km0graytext} Version 1.0}

\vfill

\begin{tikzpicture}
\draw[km0blue!40, line width=0.6pt] (-7,0) -- (7,0);
\end{tikzpicture}
\vspace{0.5cm}

\begin{tabularx}{0.85\linewidth}{@{}X X@{}}
\textbf{Réalisé par~:} & \textbf{Encadré par~:} \\[0.15cm]
\nomEtudiant & \nomEncadrant \\
\end{tabularx}

\vspace{1.2cm}
{\small\color{km0graytext} \villeEcole\ --- Année universitaire \anneeUniv}

\vspace{0.5cm}
\end{titlepage}

\tableofcontents
\clearpage

% ================================================================
\chapter{Contexte et présentation du projet}
% ================================================================

\section{Présentation}

\textbf{KM0} est une application web de gestion destinée à un garage automobile marocain
travaillant à la fois pour des \textbf{particuliers} et pour des \textbf{institutions étatiques}.
Elle a pour objectif de remplacer un suivi actuellement informel (papier, fichiers Excel,
mémoire des employés) par un outil centralisé, traçable et sécurisé, couvrant l'intégralité
du cycle de vie d'un véhicule pris en charge par le garage : de son entrée jusqu'à sa sortie.

Le garage traite quatre grandes familles d'interventions :

\begin{table}[h!]
\centering
\rowcolors{2}{km0gray}{white}
\begin{tabularx}{\linewidth}{@{}l>{\raggedright\arraybackslash}X@{}}
\toprule
\rowcolor{km0blue}
\textcolor{white}{\textbf{Type d'intervention}} & \textcolor{white}{\textbf{Description}} \\
\midrule
\textbf{Diagnostic} & Identification d'une panne ou d'un dysfonctionnement. \\
\textbf{Réparation} & Intervention mécanique : moteur, transmission, etc. \\
\textbf{Entretien} & Vidange, changement de filtres, révision périodique. \\
\textbf{Tôlerie / Carrosserie} & Véhicules accidentés, redressage, peinture. \\
\bottomrule
\end{tabularx}
\caption{Types d'interventions couvertes par le système}
\end{table}

Un même dossier peut cumuler plusieurs types d'intervention (par exemple un diagnostic
suivi d'une réparation).

\section{Problématique actuelle}

Le fonctionnement actuel du garage repose sur des outils non structurés, ce qui engendre
plusieurs difficultés opérationnelles :

\begin{table}[h!]
\centering
\rowcolors{2}{km0gray}{white}
\begin{tabularx}{\linewidth}{@{}>{\raggedright\arraybackslash}X>{\raggedright\arraybackslash}X@{}}
\toprule
\rowcolor{km0blue}
\textcolor{white}{\textbf{Problème constaté}} & \textcolor{white}{\textbf{Conséquence}} \\
\midrule
Suivi papier / Excel / mémoire des employés & Aucune traçabilité, perte d'information \\
Pas de contrôle formalisé à la sortie du véhicule & Véhicules sortis sans paiement ou sans autorisation \\
Historique client dispersé & Impossible de savoir ce qui a été fait sur un véhicule \\
Aucune visibilité pour le gérant & Pas de vue sur la charge de travail, le chiffre d'affaires, les délais \\
\bottomrule
\end{tabularx}
\caption{Problèmes actuels et conséquences}
\end{table}

\section{Objectifs du projet}

Le projet KM0 poursuit quatre objectifs principaux :

\begin{enumerate}[label=\arabic*., leftmargin=1.2cm]
\item \textbf{Centraliser} la gestion des dossiers véhicules, des clients et des interventions.
\item \textbf{Sécuriser la sortie} des véhicules via un \emph{Bon de Sortie} (BS) vérifiable,
      empêchant toute sortie non autorisée ou non payée.
\item \textbf{Offrir à toute l'équipe une meilleure traçabilité et organisation} du travail
      quotidien : chaque intervenant (réceptionniste, technicien, admin) sait à tout moment
      ce qui doit être fait et ce qui a déjà été fait sur un dossier.
\item \textbf{Donner au gérant une vision business exploitable} : charge de l'atelier,
      chiffre d'affaires, délais de traitement.
\end{enumerate}

\section{Présentation de l'organisme d'accueil}

\begin{table}[h!]
\centering
\rowcolors{2}{km0gray}{white}
\begin{tabularx}{\linewidth}{@{}l>{\raggedright\arraybackslash}X@{}}
\toprule
\rowcolor{km0blue}
\textcolor{white}{\textbf{Champ}} & \textcolor{white}{\textbf{Valeur}} \\
\midrule
Gérant & Marouane Essaadouni \\
Activité & Mécanicien réparateur \\
Adresse & 24, Allée des Tamaris, Aïn Sebaâ, Casablanca, Maroc \\
Registre de Commerce (RC) & 397421 \\
Tribunal de commerce & Casablanca \\
ICE & 002027113000005 \\
Capital social & 100 000 DHS \\
\bottomrule
\end{tabularx}
\caption{Fiche signalétique de l'entreprise d'accueil}
\end{table}

% ================================================================
\chapter{Flux opérationnel}
% ================================================================

Cette section décrit la logique de bout en bout du traitement d'un véhicule, du dépôt
par le client jusqu'à sa sortie. C'est cette logique qui structure l'ensemble des
fonctionnalités du système.

\section{Vue d'ensemble}

\begin{figure}[h!]
\centering
\begin{adjustbox}{max width=\linewidth}
\begin{tikzpicture}[
  step/.style={rectangle, rounded corners=3pt, draw=km0blue, thick, fill=km0blue!7,
               minimum width=2.9cm, minimum height=1.1cm, align=center, font=\scriptsize},
  final/.style={rectangle, rounded corners=3pt, draw=km0green, thick, fill=km0green!12,
               minimum width=2.9cm, minimum height=1.1cm, align=center, font=\scriptsize},
  arr/.style={-Stealth, thick, km0graytext}
]
\node[step] (arrivee) {Arrivée du véhicule};
\node[step, right=1cm of arrivee] (dossier) {Création du dossier\\(réceptionniste)};
\node[step, right=1cm of dossier] (fiche) {Fiche d'intervention\\générée};
\node[step, right=1cm of fiche] (charge) {Prise en charge\\(technicien)};
\node[step, right=1cm of charge] (avance) {Avancement\\(Kanban)};
\node[final, right=1cm of avance] (sortie) {BS $\rightarrow$ Sortie};

\draw[arr] (arrivee) -- (dossier);
\draw[arr] (dossier) -- (fiche);
\draw[arr] (fiche) -- (charge);
\draw[arr] (charge) -- (avance);
\draw[arr] (avance) -- (sortie);
\end{tikzpicture}
\end{adjustbox}
\caption{Logique de bout en bout du traitement d'un véhicule}
\end{figure}

\section{Prise en charge et fiche d'intervention}

Dès la création du dossier, le système génère automatiquement une \textbf{fiche
d'intervention} : un document résumé, imprimable, destiné à être placé physiquement
dans le véhicule (pare-brise ou tableau de bord). Son rôle est de donner à tout
technicien qui prend le véhicule en charge une vision immédiate de ce qui est attendu,
sans avoir à consulter le système pour les informations de base.

\begin{itemize}[leftmargin=1cm]
\item \textbf{Génération} : automatique, dès que le dossier passe au statut \texttt{ENTRE}.
\item \textbf{Contenu} : numéro de dossier, immatriculation, marque/modèle, client,
      motif déclaré, type(s) d'intervention demandée(s), date d'entrée.
\item \textbf{Nature} : document purement informatif, sans QR code ni valeur d'autorisation
      --- à la différence du Bon de Sortie, il ne conditionne aucune action.
\item \textbf{Source des données} : aucune nouvelle entité n'est nécessaire, la fiche est
      générée directement à partir des champs existants du modèle \texttt{Dossier}.
\end{itemize}

Une fois le véhicule pris en charge, son avancement (diagnostic, devis, réparation,
prêt) est visible en temps réel dans la vue Kanban de l'atelier et tracé dans
l'historique du dossier (voir chapitre \ref{ch:regles}).

% ================================================================
\chapter{Périmètre du projet}
% ================================================================

\section{Fonctionnalités couvertes}

Le projet couvre l'ensemble du cycle de vie d'un dossier :

\begin{itemize}[leftmargin=1cm]
\item Authentification et gestion des rôles
\item Gestion des clients et des véhicules
\item Gestion des dossiers et vue Kanban de l'atelier
\item Génération automatique d'une fiche d'intervention à la prise en charge du véhicule
\item Génération du Bon de Sortie (PDF + QR code)
\item Écran de vérification de sortie (scan QR)
\item Facturation
\item Tableau de bord de synthèse
\item Historique complet par véhicule
\end{itemize}

% ================================================================
\chapter{Architecture technique}
% ================================================================

\section{Stack technique}

\begin{table}[h!]
\centering
\rowcolors{2}{km0gray}{white}
\begin{tabularx}{\linewidth}{@{}l>{\raggedright\arraybackslash}X@{}}
\toprule
\rowcolor{km0blue}
\textcolor{white}{\textbf{Composant}} & \textcolor{white}{\textbf{Technologie}} \\
\midrule
Framework applicatif & Next.js 15 (App Router) + TypeScript (mode strict) \\
Base de données & PostgreSQL 16 \\
ORM & Prisma 6 \\
Interface utilisateur & Tailwind CSS + shadcn/ui \\
Authentification & Session personnalisée (cookies \texttt{httpOnly} + \texttt{bcrypt}) \\
Génération PDF & \texttt{@react-pdf/renderer} \\
QR Code & \texttt{qrcode} (génération) + \texttt{html5-qrcode} (lecture navigateur) \\
Validation des données & Zod \\
Gestion des dates & \texttt{date-fns} (locale \texttt{fr}) \\
Tests & Vitest \\
\bottomrule
\end{tabularx}
\caption{Stack technique retenue}
\end{table}

\section{Contraintes techniques}

\begin{itemize}[leftmargin=1cm]
\item \textbf{Aucune dépendance cloud tierce} pour l'authentification (pas de Clerk, Auth0,
      Supabase Auth) : l'application doit rester autonome.
\item Les \textbf{Server Actions} sont privilégiées par rapport aux routes API. Les seules
      exceptions sont l'endpoint de vérification du Bon de Sortie, qui doit rester accessible
      via une route API classique pour être appelable par le scanner, et les endpoints de
      génération de documents PDF (fiche d'intervention, Bon de Sortie), qui nécessitent une
      route API pour renvoyer un flux binaire téléchargeable.
\item Aucun type \texttt{any} n'est toléré en TypeScript ; la commande \texttt{tsc --noEmit}
      doit systématiquement passer sans erreur.
\end{itemize}

\section{Schéma d'architecture}

\begin{figure}[h!]
\centering
\begin{adjustbox}{max width=\linewidth}
\begin{tikzpicture}[
  box/.style={rectangle, rounded corners=3pt, draw=km0blue, thick, fill=km0blue!6,
              minimum width=3.4cm, minimum height=1.1cm, align=center, font=\small},
  boxalt/.style={rectangle, rounded corners=3pt, draw=km0orange, thick, fill=km0orange!8,
              minimum width=3.4cm, minimum height=1.1cm, align=center, font=\small},
  db/.style={cylinder, shape border rotate=90, aspect=0.25, draw=km0blue, thick,
             fill=km0blue!10, minimum width=2.6cm, minimum height=1.4cm, align=center, font=\small},
  arr/.style={-Stealth, thick, km0graytext}
]
\node[box] (client) {Navigateur\\(réceptionniste, technicien, admin)};
\node[boxalt, right=2.4cm of client] (vigile) {Tablette guérite\\(vigile)};

\node[box, below=1.3cm of $(client)!0.5!(vigile)$, minimum width=6.5cm] (next) {Next.js 15 --- App Router\\Server Actions};
\node[boxalt, below=0.55cm of next, minimum width=6.5cm] (api) {\texttt{/api/verification} (route API)};

\node[box, below=0.9cm of api, xshift=-2.6cm] (prisma) {Prisma ORM};
\node[box, below=0.9cm of api, xshift=2.6cm] (services) {Services métier\\(\texttt{lib/services})};

\node[db, below=1.1cm of $(prisma)!0.5!(services)$] (pg) {PostgreSQL 16};

\node[boxalt, right=2.8cm of services] (pdf) {PDF + QR\\(BS)};

\draw[arr] (client) -- (next);
\draw[arr] (vigile) -- (api);
\draw[arr] (next) -- (api);
\draw[arr] (next) -- (services);
\draw[arr] (api) -- (services);
\draw[arr] (services) -- (prisma);
\draw[arr] (prisma) -- (pg);
\draw[arr] (services) -- (pdf);
\end{tikzpicture}
\end{adjustbox}
\caption{Architecture applicative simplifiée}
\end{figure}

% ================================================================
\chapter{Modèle de données}
% ================================================================

\section{Diagramme entité-relation}

\begin{sidewaysfigure}[p]
\centering
\resizebox{\textheight}{!}{%
\begin{tikzpicture}[
  entity/.style={rectangle, draw=km0blue, thick, fill=white, rounded corners=2pt,
                 align=left, font=\footnotesize, drop shadow},
  title/.style={font=\bfseries\small\color{white}, fill=km0blue},
  rel/.style={-Stealth, thick, km0graytext, font=\scriptsize},
  every node/.style={inner sep=4pt}
]

% ---- User ----
\node[entity] (user) at (0,7) {
  \begin{tabular}{@{}l@{}}
  \multicolumn{1}{c}{\cellcolor{km0blue}\textcolor{white}{\textbf{User}}}\\
  id, nom, email\\ passwordHash, role\\ actif
  \end{tabular}
};

% ---- Client ----
\node[entity] (client) at (11,9) {
  \begin{tabular}{@{}l@{}}
  \multicolumn{1}{c}{\cellcolor{km0blue}\textcolor{white}{\textbf{Client}}}\\
  id, nom, telephone\\ email, cin, type\\ raisonSociale, ice
  \end{tabular}
};

% ---- Vehicule ----
\node[entity] (vehicule) at (11,6) {
  \begin{tabular}{@{}l@{}}
  \multicolumn{1}{c}{\cellcolor{km0blue}\textcolor{white}{\textbf{Vehicule}}}\\
  id, immatriculation\\ marque, modele, annee\\ couleur, kilometrage
  \end{tabular}
};

% ---- Dossier (hub central) ----
\node[entity, fill=km0blue!8] (dossier) at (11,2.5) {
  \begin{tabular}{@{}l@{}}
  \multicolumn{1}{c}{\cellcolor{km0blue}\textcolor{white}{\textbf{Dossier}}}\\
  id, numero\\ typesIntervention[]\\ statut, motifDeclare\\
  dateEntree, datePrevueSortie\\ dateSortieReelle\\
  montantTotal, statutPaiement
  \end{tabular}
};

% ---- BonDeSortie ----
\node[entity] (bs) at (0,3.6) {
  \begin{tabular}{@{}l@{}}
  \multicolumn{1}{c}{\cellcolor{km0blue}\textcolor{white}{\textbf{BonDeSortie}}}\\
  id, numero, qrToken\\ statut\\ dateGeneration\\ dateUtilisation\\
  derogationPaiement\\ derogationMotif
  \end{tabular}
};

% ---- Facture ----
\node[entity] (facture) at (20,4.5) {
  \begin{tabular}{@{}l@{}}
  \multicolumn{1}{c}{\cellcolor{km0blue}\textcolor{white}{\textbf{Facture}}}\\
  id, numero\\ montantHT, tauxTVA\\ montantTTC\\
  statutPaiement\\ modePaiement, datePaiement
  \end{tabular}
};

% ---- PieceUtilisee ----
\node[entity] (piece) at (20,1.6) {
  \begin{tabular}{@{}l@{}}
  \multicolumn{1}{c}{\cellcolor{km0blue}\textcolor{white}{\textbf{PieceUtilisee}}}\\
  id, nomPiece, reference\\ quantite, prixUnitaire
  \end{tabular}
};

% ---- HistoriqueStatut ----
\node[entity] (hist) at (11,-0.8) {
  \begin{tabular}{@{}l@{}}
  \multicolumn{1}{c}{\cellcolor{km0blue}\textcolor{white}{\textbf{HistoriqueStatut}}}\\
  id, statutAvant, statutApres\\ commentaire, createdAt
  \end{tabular}
};

% ---- Relations ----
\draw[rel] (client) -- (vehicule) node[midway, above, sloped]{1 -- N};
\draw[rel] (vehicule) -- (dossier) node[midway, right]{1 -- N};
\draw[rel] (dossier) -- (bs) node[midway, above, sloped]{1 -- N};
\draw[rel] (dossier) -- (facture) node[midway, above, sloped]{1 -- N};
\draw[rel] (dossier) -- (piece) node[midway, below, sloped]{1 -- N};
\draw[rel] (dossier) -- (hist) node[midway, right]{1 -- N};
\draw[rel] (user) -- (dossier) node[midway, above, sloped, xshift=-0.4cm]{1 -- N (technicien, optionnel)};
\draw[rel] (user) -- (hist) node[midway, below, sloped]{1 -- N (auteur)};
\draw[rel] (user) -- (bs) node[midway, left]{1 -- N (dérogation, optionnel)};

\end{tikzpicture}%
}
\caption{Diagramme entité-relation du modèle KM0}
\end{sidewaysfigure}

\section{Description des entités principales}

\begin{table}[h!]
\centering
\rowcolors{2}{km0gray}{white}
\begin{tabularx}{\linewidth}{@{}l>{\raggedright\arraybackslash}X@{}}
\toprule
\rowcolor{km0blue}
\textcolor{white}{\textbf{Entité}} & \textcolor{white}{\textbf{Rôle}} \\
\midrule
\textbf{User} & Compte utilisateur interne (réceptionniste, technicien, admin, vigile). \\
\textbf{Client} & Particulier, entreprise ou institution étatique propriétaire de véhicules. \\
\textbf{Vehicule} & Véhicule rattaché à un client, identifié par son immatriculation. \\
\textbf{Dossier} & Intervention en cours ou passée sur un véhicule ; entité centrale du système. \\
\textbf{BonDeSortie} & Autorisation de sortie du véhicule, matérialisée par un QR code unique. \\
\textbf{Facture} & Document de facturation rattaché à un dossier. \\
\textbf{PieceUtilisee} & Pièce détachée consommée dans le cadre d'un dossier. \\
\textbf{HistoriqueStatut} & Trace d'audit de chaque changement de statut d'un dossier. \\
\bottomrule
\end{tabularx}
\caption{Rôle de chaque entité du modèle}
\end{table}

\section{Énumérations}

\begin{table}[h!]
\centering
\rowcolors{2}{km0gray}{white}
\begin{tabularx}{\linewidth}{@{}l>{\raggedright\arraybackslash}X@{}}
\toprule
\rowcolor{km0blue}
\textcolor{white}{\textbf{Énumération}} & \textcolor{white}{\textbf{Valeurs possibles}} \\
\midrule
\texttt{Role} & \texttt{RECEPTIONNISTE}, \texttt{TECHNICIEN}, \texttt{ADMIN}, \texttt{VIGILE} \\
\texttt{TypeIntervention} & \texttt{DIAGNOSTIC}, \texttt{REPARATION}, \texttt{ENTRETIEN}, \texttt{TOLERIE} \\
\texttt{StatutDossier} & \texttt{ENTRE}, \texttt{EN\_DIAGNOSTIC}, \texttt{DEVIS\_ENVOYE}, \texttt{DEVIS\_ACCEPTE}, \texttt{DEVIS\_REFUSE}, \texttt{EN\_COURS}, \texttt{PRET}, \texttt{SORTI}, \texttt{ANNULE} \\
\texttt{StatutPaiement} & \texttt{NON\_SOLDE}, \texttt{PARTIEL}, \texttt{SOLDE} \\
\texttt{ModePaiement} & \texttt{ESPECES}, \texttt{CHEQUE}, \texttt{VIREMENT}, \texttt{CARTE} \\
\texttt{StatutBS} & \texttt{VALIDE}, \texttt{UTILISE}, \texttt{ANNULE} \\
\texttt{TypeClient} & \texttt{PARTICULIER}, \texttt{ETAT}, \texttt{ENTREPRISE} \\
\bottomrule
\end{tabularx}
\caption{Énumérations du modèle de données}
\end{table}

% ================================================================
\chapter{Rôles et permissions}
% ================================================================

Le contrôle d'accès repose sur quatre rôles distincts. Il est appliqué à deux niveaux :
via le middleware Next.js pour le routage, et systématiquement à l'intérieur de chaque
Server Action (défense en profondeur).

\begin{table}[h!]
\centering
\rowcolors{2}{km0gray}{white}
\begin{tabularx}{\linewidth}{@{}l>{\raggedright\arraybackslash}X@{}}
\toprule
\rowcolor{km0blue}
\textcolor{white}{\textbf{Rôle}} & \textcolor{white}{\textbf{Permissions}} \\
\midrule
\textbf{RECEPTIONNISTE} & Créer et éditer clients, véhicules, dossiers. Voir tous les dossiers.
Encaisser les paiements. \textbf{Ne peut pas} générer de Bon de Sortie. \\
\textbf{TECHNICIEN} & Voir les dossiers qui lui sont assignés. Changer le statut
(\texttt{EN\_DIAGNOSTIC} $\rightarrow$ \texttt{EN\_COURS} $\rightarrow$ \texttt{PRET}).
Ajouter des observations et des pièces utilisées. \\
\textbf{ADMIN} & Accès total. Génération du Bon de Sortie. Dérogations de paiement.
Consultation des statistiques. Gestion des utilisateurs. Régénération d'un BS perdu. \\
\textbf{VIGILE} & Accès \textbf{exclusif} à l'écran \texttt{/verification}. Redirection forcée
depuis toute autre route. Aucune donnée financière visible. \\
\bottomrule
\end{tabularx}
\caption{Matrice des rôles et permissions}
\end{table}

% ================================================================
\chapter{Règles métier critiques}
\label{ch:regles}
% ================================================================

Les règles suivantes sont \textbf{non négociables} : elles conditionnent l'intégrité et la
sécurité du processus de sortie des véhicules.

\begin{regle}{1}{Statut \texttt{SORTI} automatique uniquement}
Le statut \texttt{SORTI} n'est posé \textbf{que} via le scan du Bon de Sortie sur l'écran
\texttt{/verification}. Aucune interface ne permet de le poser manuellement.
\end{regle}

\begin{regle}{2}{Génération du BS conditionnée}
Le bouton de génération n'est visible que si le dossier est au statut \texttt{PRET}.
Si le statut de paiement n'est pas \texttt{SOLDE}, la génération est bloquée, sauf
dérogation d'un Admin (motif obligatoire, 10 caractères minimum, tracé dans le système).
\end{regle}

\begin{regle}{3}{Usage unique du BS}
Au moment du scan, le passage au statut \texttt{UTILISE} est effectué de manière atomique
(transaction Prisma). Un même QR code ne peut jamais autoriser deux sorties, même en cas
de double scan simultané.
\end{regle}

\begin{regle}{4}{Régénération d'un BS perdu}
En cas de perte, seul un Admin peut générer un nouveau Bon de Sortie ; l'ancien passe au
statut \texttt{ANNULE} dans la même transaction.
\end{regle}

\begin{regle}{5}{Multi-intervention, BS unique}
Un dossier peut cumuler plusieurs types d'intervention (\texttt{typesIntervention}), mais
ne génère jamais qu'un seul Bon de Sortie.
\end{regle}

\begin{regle}{6}{Audit trail obligatoire}
Tout changement de statut est écrit dans \texttt{HistoriqueStatut} avec l'utilisateur
responsable. Ce mécanisme est implémenté au niveau du service métier et n'est pas
contournable depuis l'interface.
\end{regle}

\begin{regle}{7}{Numérotation séquentielle}
Les identifiants \texttt{DOS-2026-XXXX}, \texttt{BS-2026-XXXX} et \texttt{FACT-2026-XXXX}
sont générés séquentiellement par année, sans trou, via une transaction dédiée afin
d'éviter toute collision.
\end{regle}

\section{Machine à états du dossier}

\begin{figure}[h!]
\centering
\begin{adjustbox}{max width=\linewidth}
\begin{tikzpicture}[
  st/.style={rectangle, rounded corners=3pt, draw=km0blue, thick, fill=km0blue!7,
             minimum width=2.6cm, minimum height=0.9cm, align=center, font=\scriptsize},
  final/.style={rectangle, rounded corners=3pt, draw=km0green, thick, fill=km0green!12,
             minimum width=2.6cm, minimum height=0.9cm, align=center, font=\scriptsize},
  cancel/.style={rectangle, rounded corners=3pt, draw=km0red, thick, fill=km0red!10,
             minimum width=2.6cm, minimum height=0.9cm, align=center, font=\scriptsize},
  tr/.style={-Stealth, thick, km0graytext, font=\tiny},
  trc/.style={-Stealth, thick, km0red!70, dashed, font=\tiny}
]
\node[st] (entre) {ENTRE};
\node[st, right=1.1cm of entre] (diag) {EN\_DIAGNOSTIC};
\node[st, right=1.1cm of diag] (devis) {DEVIS\_ENVOYE};
\node[st, above right=0.7cm and 1.3cm of devis] (accepte) {DEVIS\_ACCEPTE};
\node[st, below right=0.7cm and 1.3cm of devis] (refuse) {DEVIS\_REFUSE};
\node[st, right=1.6cm of accepte] (cours) {EN\_COURS};
\node[st, right=1.1cm of cours] (pret) {PRET};
\node[final, right=1.4cm of pret] (sorti) {SORTI};
\node[cancel, below=1.6cm of devis] (annule) {ANNULE};

\draw[tr] (entre) -- (diag);
\draw[tr] (diag) -- (devis);
\draw[tr] (devis) -- (accepte);
\draw[tr] (devis) -- (refuse);
\draw[tr] (accepte) -- (cours);
\draw[tr] (cours) -- (pret);
\draw[tr] (pret) -- node[above, font=\tiny\bfseries, text=km0green]{scan BS} (sorti);

\draw[trc] (diag) -- (annule);
\draw[trc] (devis) -- (annule);
\draw[trc] (refuse) -- (annule);
\draw[trc] (cours) -- (annule);
\end{tikzpicture}
\end{adjustbox}
\caption{Transitions de statut autorisées pour un dossier (traits pleins = flux normal,
traits pointillés rouges = annulation)}
\end{figure}

Toute transition de statut, y compris vers \texttt{ANNULE}, déclenche une écriture dans
\texttt{HistoriqueStatut} (règle 6). Seule la transition vers \texttt{SORTI} échappe aux
interfaces classiques de changement de statut : elle est réservée à l'endpoint de
vérification (règle 1).

% ================================================================
\chapter{Écran de vérification}
% ================================================================

\section{Contexte d'usage}

L'écran \texttt{/verification} est l'écran le plus critique du système. Il est utilisé par
un agent de sécurité \textbf{non technique}, sur tablette, dans une guérite. Son
ergonomie doit donc privilégier la lisibilité et la simplicité absolue :

\begin{itemize}[leftmargin=1cm]
\item Champ de saisie en auto-focus, combiné à un scanner QR par caméra ou à une saisie manuelle du numéro.
\item Typographie très large, lisible à 2 mètres de distance.
\item Zéro jargon technique.
\item Retour automatique à l'état de scan après 8 secondes.
\end{itemize}

\section{Logique de vérification}

À chaque scan, le serveur exécute les vérifications suivantes, dans l'ordre :

\begin{enumerate}[leftmargin=1.2cm]
\item Le Bon de Sortie existe-t-il ?
\item Son statut est-il \texttt{VALIDE} ?
\item Le dossier associé est-il au statut \texttt{PRET} ?
\item Le paiement est-il soldé, ou une dérogation a-t-elle été accordée ?
\end{enumerate}

\begin{table}[h!]
\centering
\rowcolors{2}{km0gray}{white}
\begin{tabularx}{\linewidth}{@{}l>{\raggedright\arraybackslash}X@{}}
\toprule
\rowcolor{km0blue}
\textcolor{white}{\textbf{État}} & \textcolor{white}{\textbf{Affichage}} \\
\midrule
\etat{VERT}{km0green} & « SORTIE AUTORISÉE » --- immatriculation en très grand, marque/modèle, client. \\
\etat{ORANGE}{km0orange} & « SORTIE AUTORISÉE --- DÉROGATION » --- « Non payé, autorisé par [Admin] ». \\
\etat{ROUGE}{km0red} & « SORTIE REFUSÉE » avec raison précise : facture non réglée, bon déjà
utilisé (avec date), bon inconnu, ou véhicule non prêt. \\
\bottomrule
\end{tabularx}
\caption{États d'affichage de l'écran de vérification}
\end{table}

\section{Diagramme de séquence --- génération et vérification du BS}

\begin{figure}[h!]
\centering
\begin{adjustbox}{max width=\linewidth}
\begin{tikzpicture}[
  actor/.style={rectangle, draw=km0blue, thick, fill=km0blue!8, rounded corners=2pt,
                minimum width=2.6cm, minimum height=0.8cm, align=center, font=\scriptsize\bfseries},
  lifeline/.style={dashed, km0graytext},
  msg/.style={-Stealth, thick, km0blue, font=\scriptsize},
  msgret/.style={-Stealth, thick, km0graytext, dashed, font=\scriptsize},
  note/.style={rectangle, draw=km0orange, fill=km0orange!8, rounded corners=2pt, font=\scriptsize, align=left, inner sep=4pt}
]
% Acteurs
\node[actor] (admin) at (0,0) {Admin};
\node[actor] (app) at (4.2,0) {Application\\(bs.service)};
\node[actor] (bdd) at (8.4,0) {Base de\\données};
\node[actor] (vigile) at (12.6,0) {Vigile};

% Lifelines
\draw[lifeline] (admin.south) -- ++(0,-11.5);
\draw[lifeline] (app.south) -- ++(0,-11.5);
\draw[lifeline] (bdd.south) -- ++(0,-11.5);
\draw[lifeline] (vigile.south) -- ++(0,-11.5);

% Phase 1 : génération
\draw[msg] ($(admin.south)+(0,-1)$) -- node[above]{demande génération BS (dossier PRET)} ($(app.south)+(0,-1)$);
\draw[msg] ($(app.south)+(0,-1.9)$) -- node[above]{vérifie statutPaiement} ($(bdd.south)+(0,-1.9)$);
\draw[msgret] ($(bdd.south)+(0,-2.7)$) -- node[above]{SOLDE (ou dérogation)} ($(app.south)+(0,-2.7)$);
\draw[msg] ($(app.south)+(0,-3.6)$) -- node[above]{crée BonDeSortie (qrToken, VALIDE)} ($(bdd.south)+(0,-3.6)$);
\draw[msgret] ($(app.south)+(0,-4.6)$) -- node[above]{PDF A5 + QR code} ($(admin.south)+(0,-4.6)$);

\node[note, align=center] at (4.2,-5.4) {\textit{-- temps écoulé : véhicule prêt à récupérer --}};

% Phase 2 : vérification
\draw[msg] ($(vigile.south)+(0,-6.4)$) -- node[above]{scan QR (qrToken)} ($(app.south)+(0,-6.4)$);

\node[note] at (7,-7.6) {\textbf{Transaction atomique}\\BS existe ? statut VALIDE ?\\dossier PRET ? paiement soldé ?};

\draw[msg] ($(app.south)+(0,-9)$) -- node[above]{lecture + vérifications} ($(bdd.south)+(0,-9)$);
\draw[msg] ($(app.south)+(0,-9.8)$) -- node[above]{BS $\rightarrow$ UTILISE, Dossier $\rightarrow$ SORTI} ($(bdd.south)+(0,-9.8)$);
\draw[msgret] ($(app.south)+(0,-10.8)$) -- node[above]{VERT / ORANGE / ROUGE} ($(vigile.south)+(0,-10.8)$);

\end{tikzpicture}
\end{adjustbox}
\caption{Séquence complète : génération du BS puis vérification à la sortie}
\end{figure}

% ================================================================
\chapter{Organisation du code}
% ================================================================

\section{Arborescence cible}

\begin{figure}[h!]
\centering
\begin{adjustbox}{max width=\linewidth}
\begin{tikzpicture}
\node[draw=km0blue, fill=km0gray, rounded corners=3pt, inner sep=10pt, align=left, font=\ttfamily\footnotesize]{%
\begin{tabular}{@{}l@{}}
app/ \\
\quad (auth)/login/ \\
\quad (app)/dashboard/ \\
\quad (app)/atelier/ \hfill\textcolor{km0graytext}{\% vue Kanban} \\
\quad (app)/dossiers/, dossiers/new/, dossiers/[id]/ \\
\quad (app)/clients/ \\
\quad (app)/vehicules/ \\
\quad (app)/factures/ \\
\quad (app)/admin/ \\
\quad verification/ \hfill\textcolor{km0graytext}{\% hors layout app, accès vigile} \\
\quad api/verification/route.ts \hfill\textcolor{km0graytext}{\% endpoint scan (POST)} \\
\quad api/dossiers/[id]/fiche/route.ts \hfill\textcolor{km0graytext}{\% PDF fiche d'intervention} \\
lib/ \\
\quad auth/ \hfill\textcolor{km0graytext}{\% session, hashing, guards} \\
\quad services/ \hfill\textcolor{km0graytext}{\% logique métier} \\
\quad validations/ \hfill\textcolor{km0graytext}{\% schémas Zod} \\
\quad utils/ \\
components/ \\
\quad ui/ \hfill\textcolor{km0graytext}{\% shadcn} \\
\quad features/ \hfill\textcolor{km0graytext}{\% composants métier} \\
prisma/ \\
\end{tabular}
};
\end{tikzpicture}
\end{adjustbox}
\caption{Arborescence de dossiers du projet}
\end{figure}

\section{Conventions}

\begin{itemize}[leftmargin=1cm]
\item La logique métier réside exclusivement dans \texttt{lib/services/}, jamais dans les composants React.
\item Toute Server Action valide ses entrées avec Zod avant traitement.
\item Toute Server Action vérifie le rôle de l'utilisateur en première ligne.
\item Nommage en français pour le domaine métier (\texttt{dossier}, \texttt{vehicule},
      \texttt{bonDeSortie}), en anglais pour la technique (\texttt{getUser}, \texttt{formatDate}).
\item Messages utilisateur en français ; journaux techniques (\emph{logs}) en anglais.
\item Commits au format conventionnel : \texttt{feat:}, \texttt{fix:}, \texttt{chore:},
      \texttt{test:}, \texttt{docs:}.
\end{itemize}

% ================================================================
\chapter{Planning de réalisation}
% ================================================================

Le développement est organisé sur quatre semaines, à raison d'une tâche par jour ouvré.

\begin{table}[h!]
\centering
\rowcolors{2}{km0gray}{white}
\begin{tabularx}{\linewidth}{@{}l>{\raggedright\arraybackslash}X@{}}
\toprule
\rowcolor{km0blue}
\textcolor{white}{\textbf{Semaine}} & \textcolor{white}{\textbf{Contenu}} \\
\midrule
\textbf{Semaine 1 --- Fondations} & Setup projet, authentification \& rôles, module clients,
module véhicules, création de dossier. \\
\textbf{Semaine 2 --- Cœur métier} & Vue Kanban de l'atelier, machine à états et audit,
génération du Bon de Sortie, écran de vérification, facturation. \\
\textbf{Semaine 3 --- Complétude et finition} & Tableau de bord, pièces \& détail dossier,
administration, polish UI/UX, corrections. \\
\textbf{Semaine 4 --- Tests et livraison} & Tests métier, tests de parcours de bout en bout,
données de démonstration, documentation, finalisation. \\
\bottomrule
\end{tabularx}
\caption{Planning de développement (détail complet dans \texttt{TASKS.md})}
\end{table}

% ================================================================
\chapter{Annexe --- Glossaire}
% ================================================================

\begin{table}[h!]
\centering
\rowcolors{2}{km0gray}{white}
\begin{tabularx}{\linewidth}{@{}l>{\raggedright\arraybackslash}X@{}}
\toprule
\rowcolor{km0blue}
\textcolor{white}{\textbf{Terme}} & \textcolor{white}{\textbf{Définition}} \\
\midrule
\textbf{Dossier} & Fiche représentant une intervention sur un véhicule, du dépôt à la sortie. \\
\textbf{Bon de Sortie (BS)} & Document autorisant la sortie physique d'un véhicule, muni d'un QR code à usage unique. \\
\textbf{Fiche d'intervention} & Document résumé généré à l'entrée du véhicule, placé physiquement dans celui-ci, informant le technicien du travail demandé. \\
\textbf{Dérogation} & Autorisation exceptionnelle accordée par un Admin pour sortir un véhicule non soldé. \\
\textbf{Vigile} & Agent de sécurité chargé de contrôler les sorties de véhicules à la guérite. \\
\textbf{Kanban} & Vue en colonnes représentant les dossiers selon leur statut d'avancement. \\
\textbf{Institution étatique} & Client de type organisme public (véhicules de fonction, administrations). \\
\bottomrule
\end{tabularx}
\caption{Glossaire des termes métier}
\end{table}

\end{document}
